% Abhakseite „Das kann ich“ – Zone nur im Gesamt (\mitzone)
%% TODO Ich-kann-Sätze: Platzhalter wie in den Aufgabendateien
\begin{abhakseite}
\ifdefined\mitzone
\abhakgruppe{Kennst du schon}
\abhak{1}{Längeneinheiten umrechnen (mm, cm, m), gemischte Angaben angleichen}
\abhak{2}{Flächeneinheiten (cm², m²) und Umrechnungszahl 100}
\abhak{3}{Multiplizieren und Dividieren mit Dezimalzahlen (Kommazahl mal Kommazahl, zweistellig geteilt durch einstellig)}
\abhak{4}{Formel nach einer Größe umstellen (A = a · b → b = A : a)}
\abhak{5}{Wurzel ziehen bei Quadratzahlen (√49)}
\abhak{6}{Rechten Winkel erkennen und einzeichnen (Geodreieck)}
\abhak{7}{Kreisfläche (π · r²)}
\abhak{8}{Flächeneinheiten (cm², m²) und Umrechnungszahl 100 – Fehler finden}
\abhak{9}{Flächeneinheiten (cm², m²) und Umrechnungszahl 100 – selbst rechnen}
\fi
\abhakgruppe{Einheit 1 · Rechteck, Quadrat, Umfang}
\abhak{10}{Gegeben – gesucht}
\abhak{11}{Rechteck}
\abhak{12}{aus Rechtecken zusammengesetzte Fläche (Summe)}
\abhak{13}{Einheit wechseln vor dem Rechnen (cm und m gemischt)}
\abhak{14}{Rechteck – Fehler finden, Begründen, Anwendung, Darstellungswechsel}
\abhakgruppe{Einheit 2 · Parallelogramm}
\abhak{15}{Parallelogramm}
\abhak{16}{Umfang aus den Seiten}
\abhak{17}{Parallelogramm – Fehler finden, Begründen, Anwendung}
\abhakgruppe{Einheit 3 · Dreieck}
\abhak{18}{Dreieck}
\abhak{19}{Höhe zur passenden Grundseite wählen (drei Höhen)}
\abhak{20}{Umfang}
\abhak{21}{Dreieck – Fehler finden, Begründen, Anwendung, Darstellungswechsel}
\abhakgruppe{Einheit 4 · Trapez, Drachenviereck, Raute}
\abhak{22}{Trapez}
\abhak{23}{Umfang mit Schenkel aus Pythagoras oder Sinus (Vorrat)}
\abhak{24}{Diagonale aus A}
\abhak{25}{Trapez – Fehler finden, Begründen, Anwendung}
\abhakgruppe{Einheit 5 · Zusammengesetzte Figuren}
\abhak{26}{Zusammengesetzte Figuren}
\abhak{27}{Sachaufgabe mit Entscheidung (reicht die Platte?)}
\abhak{28}{Verschnitt in Prozent (Vorrat, Niveau III)}
\abhak{29}{Zusammengesetzte Figuren – Fehler finden, Begründen, Anwendung}
\end{abhakseite}
